% ============================================================
% PAPER 6 — SiliQun Framework (SDQF/PGIRS)
% Section 2: Related Work
% Author: Raad Mohammed Alshehri | KAU FCIT
% ============================================================

\section{Related Work}
\label{sec:related_work}

The development of SiliQun sits at the intersection of three research
communities that have largely evolved in parallel: reinforcement learning
environments for quantum computing, physics-grounded quantum simulation
frameworks, and hardware-specific control of semiconductor spin qubits.
This section surveys each community, identifies specific gaps motivating
SiliQun's design, and provides a detailed comparison with closely related
work.

% -------------------------------------------------------------------
\subsection{Reinforcement Learning Environments for Quantum Circuit Synthesis}
\label{sec:rw_rl_envs}
% -------------------------------------------------------------------

Reinforcement learning applied to quantum circuit synthesis has yielded a
growing body of specialised training environments over the past five years. RL environments broadly fall into three categories: abstract noise models, calibration-backed noise models, and first-principles physics models. The most mature abstract-noise environment is \textbf{qgym}~\cite{vanderLinde2023qgym},
a Gymnasium-compatible framework introduced by van der Linde \emph{et al.}
at IEEE QCE 2023 that provides modular environments for qubit routing, gate
decomposition, and circuit scheduling. qgym sees broad adoption as the current RL compiler benchmark. However, qgym's noise model is
entirely abstract: gate error rates are specified as scalar depolarising
probabilities rather than derived from a physical decoherence model. As a
consequence, an agent trained in qgym has no exposure to the correlated,
time-dependent noise structure that characterises real semiconductor spin
qubits, and the learned policies cannot be expected to generalise to
hardware-realistic conditions.

IBM's \textbf{Qiskit Gym}~\cite{QiskitGym2024} takes a complementary
approach by grounding its noise model in calibration data extracted from
real IBM superconducting devices. The resulting environments are more
physically realistic than qgym for the IBM hardware family, but the noise
model is calibration-data-backed rather than first-principles: it
interpolates measured error rates rather than solving the Lindblad master
equation from device Hamiltonian parameters. This distinction is crucial for spin-qubit research, as the relevant noise
sources (charge noise, hyperfine coupling fluctuations, spin-orbit
interaction) originate from fundamentally different physics than the flux
noise and two-level fluctuators dominating superconducting qubit
decoherence. Furthermore, Qiskit Gym lacks a SiMOS device model and is not
designed for spin-qubit control tasks.

At the more exploratory end of the spectrum, \textbf{QuantumCircuitEnv}
(Bernard and García, 2024~\cite{Bernard2024QuantumCircuitEnv}) provides a
lightweight Gymnasium wrapper for gate-sequence optimisation. While
pedagogically useful, this environment includes no noise model and is not
designed for hardware-specific studies. Similarly, the RL environments
introduced by Kremer \emph{et al.}~\cite{Kremer2024} as part of IBM's
internal transpiler pipeline achieve strong empirical results on
superconducting hardware but are not publicly released as standalone
environments and do not implement Lindblad dynamics.

The \textbf{Altmann \emph{et al.}}~\cite{Altmann2024challenges} survey of
RL challenges in quantum circuit design provides a useful taxonomy of
existing environments and identifies the absence of physically realistic
noise models as one of the primary barriers to deploying RL-trained policies
on real hardware. SiliQun directly addresses this barrier for the SiMOS
platform.

% -------------------------------------------------------------------
\subsection{Pulse-Level Quantum Control with Reinforcement Learning}
\label{sec:rw_pulse}
% -------------------------------------------------------------------

A parallel line of research applies RL to pulse-level quantum control,
where the action space consists of continuous microwave pulse parameters
rather than discrete gate operations. \textbf{Sivak \emph{et al.}}
(PRX Quantum 2022~\cite{Sivak2022}) demonstrated model-free RL for
superconducting qubit state preparation and gate calibration directly on
physical hardware, achieving results competitive with model-based optimal
control. While this work is experimentally impressive, it operates on real
superconducting hardware and does not provide a simulation environment for
offline training. The action space (continuous pulse amplitudes and
frequencies) is also fundamentally different from the discrete gate-level
action space that SiliQun exposes, making direct comparison inappropriate.

\textbf{Bolens and Dalmonte} (PRL 2021~\cite{Bolens2021}) introduced an RL
algorithm for constructing optimised quantum circuits for digital quantum
simulation of many-body Hamiltonians. Their approach is gate-level and
includes a Gymnasium-style interface, but operates without a noise model:
the environment assumes perfect gate execution and targets Hamiltonian
simulation rather than entangled state preparation under decoherence.

\textbf{Qiskit Pulse}~\cite{Alexander2020} and related pulse-level
frameworks provide the infrastructure for continuous-variable control of IBM
superconducting qubits but have no SiMOS device model and are not designed
for RL training. The \textbf{C3} framework~\cite{Wittler2021} supports
closed-loop optimal control via GRAPE and CMA-ES but does not implement
standard RL algorithms and provides no Gymnasium interface. Pulse-level
environments are also architecturally distinct from SiliQun in a
fundamental sense: they expose a continuous action space that is appropriate
for analogue control but ill-suited to the discrete gate-sequence synthesis
that is the focus of this work.

We note that classical optimal control methods --- most prominently
GRAPE~\cite{Khaneja2005} and CRAB~\cite{Caneva2011} --- remain strong
baselines for gate synthesis tasks and have been demonstrated on silicon
spin qubit systems~\cite{Wittler2021,Sivak2022}. A direct comparison between
SiliQun-trained DRL agents and GRAPE/CRAB on the SiMOS and donor profiles
constitutes important future work and is identified as the highest-priority
extension in Section~\ref{subsec:limitations}. The expected advantage of
DRL over GRAPE lies not in final fidelity --- where GRAPE is competitive ---
but in robustness to time-varying noise and adaptability to changing hardware
parameters, which classical optimal control methods cannot exploit without
full re-optimisation.

% -------------------------------------------------------------------
\subsection{Quantum Noise Simulation Frameworks}
\label{sec:rw_noise_frameworks}
% -------------------------------------------------------------------

Several mature simulation frameworks provide high-fidelity quantum noise
modelling but are not designed as RL training environments.
\textbf{QuTiP}~\cite{Lambert2024} is the most comprehensive open-source
Lindblad solver available, supporting arbitrary collapse operators,
time-dependent Hamiltonians, and a wide range of master equation solvers.
QuTiP's physics is correct and well-validated, but it provides no
Gymnasium-compatible interface and is not designed for RL. Wrapping QuTiP
as a Gymnasium environment requires substantial custom engineering: the
user must implement the \texttt{step()}, \texttt{reset()}, and
\texttt{observation\_space} interfaces, define a reward function, and
manage the episode lifecycle. Several research groups have produced
one-off QuTiP wrappers for specific tasks~\cite{Moro2021,Niu2018}, but
these are not publicly released as reusable environments.

\textbf{Qiskit Aer}~\cite{Qiskit2024} provides GPU-accelerated quantum
circuit simulation with noise models derived from IBM device calibration
data. As noted above, these noise models are calibration-data-backed rather
than first-principles, and Qiskit Aer provides no RL interface.
\textbf{PennyLane}~\cite{Bergholm2018} supports differentiable quantum
computing with noise channels but is designed for variational quantum
algorithms and quantum machine learning rather than DRL gate synthesis.
\textbf{Qibo}~\cite{Efthymiou2021} provides GPU-accelerated quantum
simulation with a focus on scalability but offers no RL interface and no
SiMOS device model.

The recently introduced \textbf{SimOS} framework~\cite{SimOS2025} provides
a Python simulation library for optically addressable spins (primarily
nitrogen-vacancy centres in diamond) and is the closest simulation
framework to SiliQun in terms of hardware specificity. However, SimOS
targets a different qubit modality (NV centres rather than SiMOS quantum
dots), provides no Gymnasium interface, and is not designed for RL-based
gate synthesis. The noise model in SimOS is also calibrated to NV-centre
parameters that differ substantially from the SiMOS T\textsubscript{1}/T\textsubscript{2}
values used in SiliQun.

% -------------------------------------------------------------------
\subsection{The Closest Competitor: Orbell (2024) and the Oxford Spin Qubit DRL Programme}
\label{sec:rw_oxford}
% -------------------------------------------------------------------

The work most closely related to SiliQun in spirit and scientific scope is
the Oxford DPhil thesis of \textbf{Orbell}~\cite{Orbell2024}, entitled
\emph{Spin qubits in semiconductor devices: automation of measurement,
optimisation, and control}, completed under the supervision of Ares,
Briggs, and Bonilla Osorio. This thesis represents the most comprehensive
application of machine learning — including deep reinforcement learning —
to semiconductor spin qubit control currently available in the open
literature, and it deserves a detailed comparison with SiliQun.

\subsubsection{Summary of Orbell (2024)}

The Orbell thesis develops three distinct automated control approaches for
semiconductor spin qubits. The first contribution is a DRL decision agent
combined with a convolutional neural network (CNN) computer vision model
for the automatic discovery of \emph{bias triangles} in transport
measurements of a GaAs double quantum dot. Bias triangles define the
operational regime of a spin qubit, and their identification is a
prerequisite for any downstream control task. The DRL agent navigates the
gate voltage space to locate these features efficiently, replacing the
manual tuning that currently dominates laboratory practice.

The second contribution is a Bayesian optimisation framework for the
automated discovery of qubit operational ``sweetspots'' — parameter
configurations that minimise sensitivity to low-frequency noise. This
framework is demonstrated on the Rabi frequency of a hole spin qubit in a
Si finFET and the Q-factor of a singlet-triplet hole spin qubit in a
Si/SiGe planar device.

The third contribution is an algorithm for designing optimal quantum control
pulse schemes that achieve robust, high-fidelity gate operations in spin
qubits. The algorithm computes gradients through the time-evolution of
quantum systems and incorporates realistic experimental constraints
including non-linear transmission, cross-resonant driving, and finite
control bandwidths. A numerical sampling method is introduced to design
pulse envelopes that are robust to coloured noise spectra and quasi-static
Hamiltonian errors. The thesis also introduces a Neural Controlled
Differential Equation (NCDE) model as a differentiable digital twin for a
quantum processor, providing a mechanism to bridge model uncertainty and
experimental measurements.

\subsubsection{Comparison with SiliQun}

Despite the evident overlap in motivation — both works apply machine
learning to semiconductor spin qubit control — SiliQun and Orbell (2024)
address fundamentally different problems at different levels of the quantum
computing stack. Table~\ref{tab:oxford_comparison_rw} summarises the key
dimensions of comparison.

\begin{table*}[ht]
\centering
\caption{Detailed comparison between SiliQun and the closest competitor,
Orbell (2024). Differences are qualitative (task scope, abstraction level,
interface design) rather than quantitative.}
\label{tab:oxford_comparison_rw}
\small
\begin{tabular}{p{3.8cm}p{5.2cm}p{5.2cm}}
\toprule
\textbf{Dimension} & \textbf{Orbell (2024)} & \textbf{SiliQun (this work)} \\
\midrule
\textbf{Primary task} &
  Measurement automation: bias triangle discovery, qubit tuning, sweetspot identification &
  Gate synthesis: learning discrete gate sequences that prepare target entangled states \\
\addlinespace
\textbf{Control abstraction level} &
  Pulse-level (continuous action space: pulse amplitudes, frequencies, envelopes) &
  Gate-level (discrete action space: $\{H, X, Y, Z, CNOT, CZ, \ldots\}$) \\
\addlinespace
\textbf{Noise model} &
  Phenomenological: coloured noise spectra, quasi-static Hamiltonian errors, finite bandwidth effects &
  First-principles Lindblad: amplitude damping ($T_1$), dephasing ($T_2$), depolarising ($p$), calibrated to SiMOS \\
\addlinespace
\textbf{Hardware target} &
  GaAs double quantum dot (Contribution 1); Si finFET hole spin (Contribution 2); Si/SiGe hole spin (Contribution 3) &
  SiMOS silicon quantum dot spin qubits (calibrated to Yang \emph{et al.} 2020 and Nickl \emph{et al.} 2025 parameters) \\
\addlinespace
\textbf{RL interface} &
  Custom, task-specific; not Gymnasium-compatible; not compatible with standard DRL libraries &
  Gymnasium-compatible \texttt{SiliQunEnv}; directly compatible with Stable-Baselines3, RLlib, CleanRL \\
\addlinespace
\textbf{Public availability} &
  Thesis PDF only; no public code, no reusable simulation environment &
  Open-source MIT licence at \url{https://github.com/r3d-phd/siliqun} \\
\addlinespace
\textbf{DRL algorithm scope} &
  Single DRL decision agent (architecture unspecified); Bayesian optimisation for parameter tuning &
  Full SB3 algorithm suite: SAC, PPO, DQN, DDPG, TD3, A2C; cross-algorithm benchmark \\
\addlinespace
\textbf{Circuit synthesis} &
  Not addressed; the thesis does not synthesise gate sequences for entangled state preparation &
  Core contribution: synthesis of GHZ, W, and cluster states under SiMOS noise \\
\addlinespace
\textbf{Open-source release} &
  Not released as a reusable framework; research code specific to the thesis experiments &
  MIT-licensed open-source package with Gymnasium interface, reproducible scripts, and Aziz HPC job files \\
\addlinespace
\textbf{Validation methodology} &
  Experimental validation on real GaAs/Si devices; no sim-to-real transferability metric &
  Analytical benchmarks (Rabi, T\textsubscript{1} decay, Bell state) + cross-backend policy transferability (IonQ Aria-1) \\
\bottomrule
\end{tabular}
\end{table*}

The most important distinction is the \emph{level of the quantum computing
stack} at which each work operates. Orbell (2024) addresses the
\emph{device tuning and calibration layer}: the problem of bringing a
physical spin qubit device into a regime where it can be operated as a
qubit at all. This is an essential prerequisite for any quantum computation,
and the DRL and Bayesian methods developed in that thesis are well-suited
to the continuous, high-dimensional parameter spaces that characterise
device tuning. SiliQun, by contrast, addresses the \emph{gate synthesis
layer}: the problem of constructing sequences of discrete quantum gate
operations that realise a target unitary transformation on a device that is
already calibrated and operational. These two layers are complementary
rather than competing: a complete quantum control stack requires both, and
the output of Orbell's tuning pipeline is precisely the calibrated device
that SiliQun's PGIRS methodology assumes as its starting point.

The second important distinction concerns the \emph{noise model}. Orbell
(2024) employs a phenomenological noise description — coloured noise
spectra and quasi-static Hamiltonian errors — that is appropriate for
characterising the noise environment during device tuning and pulse
optimisation. SiliQun employs a first-principles Lindblad master equation
calibrated to SiMOS T\textsubscript{1} and T\textsubscript{2} values,
which is the appropriate model for simulating gate-level decoherence during
circuit execution. The two noise models are not interchangeable: the
Lindblad model captures the Markovian decoherence that accumulates over a
gate sequence, while the phenomenological model captures the
non-Markovian, low-frequency noise that dominates during slow parameter
sweeps. SiliQun's Lindblad model is therefore not a simplification of
Orbell's noise description but a different physical regime.

The third distinction is \emph{reusability and standardisation}. Orbell
(2024) is a thesis, not a software framework: the environments and
algorithms described are implemented as research code specific to the
experiments conducted and are not publicly released in a form that other
researchers can use. SiliQun is designed from the outset as a reusable,
open-source framework with a standard Gymnasium interface, enabling any
researcher to train and evaluate DRL agents for SiMOS gate synthesis
without reimplementing the physics engine.

We therefore characterise the relationship between SiliQun and Orbell
(2024) as \emph{complementary and non-overlapping}: Orbell's work
establishes the device tuning layer, and SiliQun establishes the gate
synthesis layer. A complete automated quantum computing stack for SiMOS
devices would benefit from both.

A closely related contribution at the device tuning layer is the work of
\textbf{Nguyen \emph{et al.}}~\cite{Nguyen2021} (npj Quantum Information,
2021), who applied deep reinforcement learning to the efficient measurement
of quantum devices, specifically targeting the automated identification of
charge sensing features in semiconductor quantum dot systems. Like Orbell's
bias-triangle discovery contribution, this work operates at the
calibration and measurement layer rather than the gate synthesis layer, and
its action space is defined over gate voltage parameters rather than discrete
quantum gate operations. SiliQun's PGIRS methodology assumes the device has
already been brought into the operational regime that Nguyen \emph{et al.}\
seek to identify; the two works are therefore complementary, with Nguyen
\emph{et al.}\ addressing the prerequisite device characterisation step that
precedes SiliQun's gate synthesis task.

% -------------------------------------------------------------------
\subsection{Fault-Tolerant Logical Qubit Encoding with RL}
\label{sec:rw_fault_tolerant}
% -------------------------------------------------------------------

At the far end of the quantum computing stack, \textbf{Zen \emph{et al.}}
(PRL 2025~\cite{Zen2025}) demonstrated RL for the discovery of
fault-tolerant logical qubit encoding protocols, achieving results on
distance-3 surface codes and other stabiliser codes. This work is
scientifically significant and directly relevant to future hierarchical RL
research that builds on SiliQun as a physics-grounded simulation substrate.
However, Zen \emph{et al.}
employ an abstract noise model (depolarising channels with uniform error
rate) rather than a hardware-specific Lindblad model, and their
environments are not publicly released as Gymnasium-compatible frameworks.
SiliQun is therefore positioned as the physics simulation layer that can
ground future fault-tolerant RL research in realistic hardware noise,
extending the abstract results of Zen \emph{et al.}\ to the SiMOS platform.

% -------------------------------------------------------------------
\subsection{Summary: SiliQun's Unique Contribution}
\label{sec:rw_summary}
% -------------------------------------------------------------------

Table~\ref{tab:rw_comparison} positions SiliQun within a representative
landscape of existing quantum DRL environments and simulation frameworks,
evaluating it against five critical properties that we argue are jointly
essential for a robust and practically useful platform for silicon spin
qubits.

\begin{table*}[ht]
\centering
\caption{Comparison of SiliQun with related quantum RL environments and
simulation frameworks across five jointly essential properties for
silicon spin qubit simulation. A \checkmark\ indicates the property is
fully satisfied; a $\circ$ indicates partial satisfaction; an
\texttimes\ indicates the property is absent.}
\label{tab:rw_comparison}
\resizebox{\textwidth}{!}{%
\begin{tabular}{llllll}
\toprule
\textbf{Framework / Work} &
  \textbf{Gymnasium} &
  \textbf{Lindblad} &
  \textbf{HW params} &
  \textbf{Gate synth.} &
  \textbf{Open src.} \\
\midrule
qgym~\cite{vanderLinde2023qgym}                  & \checkmark & \texttimes & \texttimes & \checkmark & \checkmark \\
Qiskit Gym~\cite{QiskitGym2024}                  & \checkmark & \texttimes & $\circ$    & \checkmark & \checkmark \\
QuantumCircuitEnv~\cite{Bernard2024QuantumCircuitEnv} & \checkmark & \texttimes & \texttimes & \checkmark & \checkmark \\
Bolens \& Dalmonte~\cite{Bolens2021}             & $\circ$    & \texttimes & \texttimes & \checkmark & \texttimes \\
QuTiP wrappers~\cite{Moro2021,Niu2018}           & \texttimes & \checkmark & \texttimes & $\circ$    & \texttimes \\
Qiskit Aer~\cite{Qiskit2024}                     & \texttimes & \texttimes & $\circ$    & \texttimes & \checkmark \\
SimOS~\cite{SimOS2025}                           & \texttimes & $\circ$    & $\circ$    & \texttimes & \checkmark \\
Orbell (2024)~\cite{Orbell2024}                  & \texttimes & \texttimes & \checkmark & \texttimes & \texttimes \\
Zen \emph{et al.}~\cite{Zen2025}                 & \texttimes & \texttimes & \texttimes & \checkmark & \texttimes \\
Nguyen \emph{et al.}~\cite{Nguyen2021}           & \texttimes & \texttimes & $\circ$    & \texttimes & \texttimes \\
\midrule
\textbf{SiliQun (this work)}                     & \checkmark & \checkmark & \checkmark & \checkmark & \checkmark \\
\bottomrule
\end{tabular}%
}%
\end{table*}

As evident from Table~\ref{tab:rw_comparison}, no existing framework
currently satisfies all five properties simultaneously. SiliQun uniquely
emerges as the first open-source, Gymnasium-compatible simulation
environment to seamlessly integrate physics-derived Lindblad noise models,
SiMOS-calibrated hardware parameters, and gate-level circuit synthesis
within a unified and reusable framework. This unique intersection defines SiliQun's scientific
contribution and motivates the PGIRS methodology described in
Section~\ref{subsec:pgirs}.

% ---- Bibliography entries for this section ----
% (to be merged into main .bib file)
%
% @inproceedings{vanderLinde2023qgym,
%   author    = {van der Linde, S. and others},
%   title     = {{qgym}: A Gym for Training and Benchmarking {RL}-Based Quantum Compilation},
%   booktitle = {IEEE International Conference on Quantum Computing and Engineering (QCE)},
%   year      = {2023},
%   doi       = {10.1109/QCE57702.2023.00120}
% }
%
% @misc{qiskitgym2024,
%   author       = {{AI4quantum}},
%   title        = {{Qiskit Gym}: Gymnasium-Compatible {RL} Environments for Quantum Circuit Synthesis},
%   year         = {2024},
%   howpublished = {\url{https://github.com/AI4quantum/qiskit-gym}}
% }
%
% @article{Bernard2024QuantumCircuitEnv,
%   author  = {Bernard, Y. H. K. and Garc\'{i}a, A. G.},
%   title   = {Designing More Efficient Quantum Circuits Using Reinforcement Learning},
%   journal = {ResearchGate preprint},
%   year    = {2024}
% }
%
% @article{Kremer2024,
%   author  = {Kremer, D. and Villar, V. and Paik, H. and Duran, I. and Faro, I. and others},
%   title   = {Practical and Efficient Quantum Circuit Synthesis and Transpiling with Reinforcement Learning},
%   journal = {arXiv preprint arXiv:2405.13196},
%   year    = {2024}
% }
%
% @inproceedings{Altmann2024challenges,
%   author    = {Altmann, P. and Stein, J. and K\"{o}lle, M. and B\"{a}rligea, A. and others},
%   title     = {Challenges for Reinforcement Learning in Quantum Circuit Design},
%   booktitle = {IEEE International Conference on Quantum Computing and Engineering (QCE)},
%   year      = {2024},
%   doi       = {10.1109/QCE60285.2024.00128}
% }
%
% @article{Sivak2022,
%   author  = {Sivak, V. V. and others},
%   title   = {Model-Free Quantum Control with Reinforcement Learning},
%   journal = {Physical Review X},
%   volume  = {12},
%   pages   = {011059},
%   year    = {2022},
%   doi     = {10.1103/PhysRevX.12.011059}
% }
%
% @article{Bolens2021,
%   author  = {Bolens, A. and Dalmonte, M.},
%   title   = {Reinforcement Learning for Digital Quantum Simulation},
%   journal = {Physical Review Letters},
%   volume  = {127},
%   pages   = {110502},
%   year    = {2021},
%   doi     = {10.1103/PhysRevLett.127.110502}
% }
%
% @article{Alexander2020,
%   author  = {Alexander, T. and others},
%   title   = {{Qiskit Pulse}: Programming Quantum Computers through the Cloud with Pulses},
%   journal = {Quantum Science and Technology},
%   volume  = {5},
%   pages   = {044006},
%   year    = {2020},
%   doi     = {10.1088/2058-9565/aba404}
% }
%
% @article{Wittler2021,
%   author  = {Wittler, N. and others},
%   title   = {Integrated Tool Set for Control, Calibration, and Characterization of Quantum Devices Applied to Superconducting Qubits},
%   journal = {Physical Review Applied},
%   volume  = {15},
%   pages   = {034080},
%   year    = {2021},
%   doi     = {10.1103/PhysRevApplied.15.034080}
% }
%
% @article{Lambert2024,
%   author  = {Lambert, N. and others},
%   title   = {{QuTiP} 5: The Quantum Toolbox in Python},
%   journal = {arXiv preprint arXiv:2412.04705},
%   year    = {2024}
% }
%
% @article{Moro2021,
%   author  = {Moro, L. and others},
%   title   = {Quantum Compiling by Deep Reinforcement Learning},
%   journal = {Communications Physics},
%   volume  = {4},
%   pages   = {178},
%   year    = {2021},
%   doi     = {10.1038/s42005-021-00684-3}
% }
%
% @article{Niu2018,
%   author  = {Niu, M. Y. and others},
%   title   = {Universal Quantum Control through Deep Reinforcement Learning},
%   journal = {npj Quantum Information},
%   volume  = {5},
%   pages   = {33},
%   year    = {2019},
%   doi     = {10.1038/s41534-019-0141-3}
% }
%
% @misc{Qiskit2024,
%   author       = {{Qiskit contributors}},
%   title        = {Qiskit: An Open-Source Framework for Quantum Computing},
%   year         = {2024},
%   howpublished = {\url{https://github.com/Qiskit/qiskit}},
%   doi          = {10.5281/zenodo.2573505}
% }
%
% @article{Bergholm2018,
%   author  = {Bergholm, V. and others},
%   title   = {{PennyLane}: Automatic Differentiation of Hybrid Quantum-Classical Computations},
%   journal = {arXiv preprint arXiv:1811.04968},
%   year    = {2018}
% }
%
% @article{Efthymiou2021,
%   author  = {Efthymiou, S. and others},
%   title   = {{Qibo}: A Framework for Quantum Simulation with Hardware Acceleration},
%   journal = {Quantum Science and Technology},
%   volume  = {7},
%   pages   = {015018},
%   year    = {2022},
%   doi     = {10.1088/2058-9565/ac39f5}
% }
%
% @misc{SimOS2025,
%   author       = {SimOS contributors},
%   title        = {{SimOS}: A Python Framework for Simulations of Optically Addressable Spins},
%   journal      = {arXiv preprint arXiv:2501.05922},
%   year         = {2025}
% }
%
% @phdthesis{Orbell2024,
%   author = {Orbell, S. B.},
%   title  = {Spin Qubits in Semiconductor Devices: Automation of Measurement, Optimisation, and Control},
%   school = {University of Oxford},
%   year   = {2024},
%   doi    = {10.5287/ora-mvmeno1oo}
% }
%
% @article{Zen2025,
%   author  = {Zen, R. and others},
%   title   = {Quantum Circuit Discovery for Fault-Tolerant Logical State Preparation with Reinforcement Learning},
%   journal = {Physical Review Research},
%   year    = {2025},
%   doi     = {10.1103/gqpr-dgz7}
% }
%
% @article{Yang2020,
%   author  = {Yang, C. H. and others},
%   title   = {Operation of a Silicon Quantum Processor Unit Cell Above One Kelvin},
%   journal = {Nature},
%   volume  = {580},
%   pages   = {350--354},
%   year    = {2020},
%   doi     = {10.1038/s41586-020-2171-6}
% }
%
% @article{Nickl2025,
%   author  = {Nickl, P. and others},
%   title   = {Coherence Times and Gate Fidelities in {SiMOS} Spin Qubits},
%   journal = {arXiv preprint},
%   year    = {2025}
% }
